\documentclass[10pt,a4paper]{amsart}
\usepackage[margin=19mm]{geometry}
\usepackage{amsmath,amssymb,mathtools}
\usepackage[scr=rsfs,cal=euler]{mathalfa}
\usepackage{xfrac}
\usepackage[unicode,hidelinks,bookmarks=false]{hyperref}
\newcommand{\ie}{i.e.\ }
\newcommand{\cf}{cf.\ }
\newcommand{\resp}{respectively\ }
\newcommand{\Subsection}{\subsection}
\providecommand{\nobreakdash}{\nobreak\hbox{-}}
\setlength{\emergencystretch}{2em}
\multlinegap0pt
\usepackage{xcolor}
\usepackage{amssymb}
\usepackage{mathtools}
\usepackage{cancel}
\usepackage[shortlabels]{enumitem}
\newtheorem{definition}{Definition}
\newtheorem{lemma}{Lemma}
\newtheorem{theorem}{Theorem}
\usepackage{cleveref}
\crefname{definition}{Definition}{Definitions}
\crefname{lemma}{Lemma}{Lemmas}
\crefname{theorem}{Theorem}{Theorems}
\newcommand{\R}{\mathbb{R}}
\title{Steerable Neural ODEs on Homogeneous Spaces}
\author{Emma Andersdotter, Daniel Persson, Fredrik Ohlsson}
\date{Source: arXiv:2605.11133v1 (CC BY 4.0)}
\begin{document}
\maketitle

\noindent\emph{Source and licence.} Steerable Neural ODEs on Homogeneous Spaces. Version arXiv:2605.11133v1 (CC BY 4.0), distributed by arXiv under the Creative Commons Attribution 4.0 International licence, http://creativecommons.org/licenses/by/4.0/, as recorded in the arXiv licence metadata for this exact version and shown on the abstract page https://arxiv.org/abs/2605.11133v1. Full licence notice: \texttt{licenses/arxiv-2605.11133-CC-BY-4.0.txt}.\par\smallskip

\noindent\textit{Editorial scope.} Counted unit: the article's theorem thm:equivariance\_steerable\_NODE\_flow (source0.tex lines 667--669). The article's complete original proof of it is delivered with it (source0.tex lines 670--686, one \texttt{proof} environment of 17 lines). No mathematics is written, repaired or re-derived: the statement and the proof are the article's own lines, in the article's own notation, and no retained line is substituted or re-flowed.  Retained uncounted as the source-local context the proof needs: lines 427--431; lines 476--482; lines 653--655.  Nothing was omitted from any retained span, so the package carries no omission record.  No retained line contains a citation, so the wrapper carries no bibliography and no bibliography entry.  No retained line contains a figure environment, an image-inclusion command or drawing code, so no image is delivered and the package declares no asset dependency.  Editorial formatting only, in addition: an AMS \texttt{amsart} preamble that restates the article's own declarations and macros used by the retained lines, namely the environments \texttt{definition}, \texttt{lemma}, \texttt{theorem} on separate counters, together with the standard packages those lines need.  The retained environments print in this document as Definition 1, Lemma 1, Theorem 1 in order of appearance, whereas the article prints the same items with the numbering of its own full document.  The complete native source of the article as distributed in the arXiv e-print for this exact version is bundled unchanged as \texttt{sources/200341/source0.tex}.
\begin{equation}\label{eq:left_action_on_M_times_V}
    L_g(p, v)
:= \varphi^{-1}\big([L_g \sigma(p), v]\big)
= (gp, \rho(c(g, p)) v)
\end{equation}

\begin{definition}[Parallel transport flow]\label{dfn:flow_ass_bundle}
    Let $M=G/H$ be a homogeneous space, $\phi$ a vector field on $M$, $\omega$ a principal $H$-connection on $G\to G/H$, and $G\times_{\rho} V$ the associated bundle defined by the representation $\rho: H \to \mathrm{GL}(V)$. Let $\Phi$ be the flow in $M$ generated by $\phi$, and $\tilde{\Phi}$ be the unique horizontal lift of $\Phi$ with respect to $\omega$. We define the map $\Gamma_{\Phi}:\R \times (G \times_{\rho} V) \to G \times_{\rho} V$ by
    \begin{equation}
        \Gamma_{\Phi}(t,[g,v]) = [\tilde{\Phi}(t,g),v]
    \end{equation}
    for every $g \in G$ and $v \in V$.
\end{definition}

\begin{lemma}\label{lem:equiv_lift_ass}
    Let $\Phi$ be the flow on the homogeneous space $M=G/H$ generated by a vector field $\phi: M \to TM$, and let $\omega$ be a principal $H$-connection on $G$. The horizontal lift $\Gamma_{\Phi}: \R \times (G \times_{\rho} V) \to G \times_{\rho} V$ of $\Phi$ in~\cref{dfn:flow_ass_bundle}, defining parallel transport in the associated bundle $G \times_{\rho} V$, is $G$-equivariant if the vector field $\phi$ and the connection $\omega$ are both $G$-invariant.
\end{lemma}

\begin{theorem}\label{thm:equivariance_steerable_NODE_flow}
    Let $M=G/H$ be a homogeneous space, let $\phi: M \to TM$ be a vector field on $M$, and let $\omega$ be a principal $H$-connection on $G$. The steerable NODE defined by $\phi$ and $\omega$ is $G$-equivariant if the vector field $\phi$ and the connection $\omega$ are both $G$-invariant.
\end{theorem}

\begin{proof}
    The steerable NODE $\Psi: M \times V \to M \times V$ can be expressed in terms of the flow $\Gamma_{\Phi}$ and the canonical local trivialisation $\phi$ with respect to the local section $\sigma$ as
    \begin{equation}
        \Psi(p,v) = \varphi^{-1} \left( \Gamma_{\Phi} (1,\varphi(p,v)) \right)\,
    \end{equation}
    for any $p \in M$ and $v \in V$. Note that this indeed defines a diffeomorphism, since both $\varphi$ and $\Gamma_{\Phi}$ are diffeomorphisms. We also note that the left action of $G$ on $M \times V$ is given by~\eqref{eq:left_action_on_M_times_V} as,
    \begin{equation}
        L_g(p,v) = \varphi^{-1} \left( L_g ( \varphi(p,v)) \right)\,,
    \end{equation}
    for any $p \in M$ and $v \in V$. Consequently,~\Cref{lem:equiv_lift_ass} implies
    \begin{eqnarray*}
        L_g \circ \Psi &=& \left(\varphi^{-1} \circ L_g \circ \varphi \right) \circ \left( \varphi^{-1} \circ \Gamma^1_{\Phi} \circ \varphi \right) = \varphi^{-1} \circ L_g \circ \Gamma^1_{\Phi} \circ \varphi \\
        &=& \varphi^{-1} \circ \Gamma^1_{\Phi} \circ L_g \circ \varphi = \left(\varphi^{-1} \circ \Gamma^1_{\Phi} \circ \varphi \right) \circ \left( \varphi^{-1} \circ L_g \circ \varphi \right) \\
        &=& \Psi \circ L_g\,,
    \end{eqnarray*}
    where we use the notation $\Gamma^1_{\Phi}([p,v]) = \Gamma_{\Phi}(1,[p,v])$. This completes the proof.
\end{proof}

\end{document}
